\documentclass[11pt,a4paper]{article}
\usepackage[utf8]{inputenc}
\usepackage{graphicx}
\usepackage{array}
\usepackage{tabularx} 
\usepackage[table,dvipsnames]{xcolor} 
\usepackage{multirow}
\usepackage{amsmath}
\usepackage{xltabular}
\usepackage{hyperref}
\usepackage{geometry}
\usepackage{float}
\usepackage{wrapfig}
\geometry{margin=1in}

% =========================================================================
% ¡AÑADIDO 1!: CONFIGURACIÓN DE LA BIBLIOGRAFÍA (PREÁMBULO)
% =========================================================================
% Cargamos el paquete con estilo IEEE (ideal para ingeniería aeroespacial)
\usepackage[style=ieee, backend=biber]{biblatex}
% Enlazamos tu archivo .bib (asegúrate de que se llame exactamente así)
\addbibresource{bibliography.bib}

% --- Configuración de colores para las tablas ---
\definecolor{EtseNavy}{HTML}{002F6C} % Azul oscuro institucional (ejemplo)
\definecolor{LightGrayTable}{HTML}{F2F4F7} % Gris claro para filas alternas

% --- Configuración de hipervínculos ---
\hypersetup{
    colorlinks=true,
    linkcolor=EtseNavy,
    filecolor=magenta,      
    urlcolor=cyan,
    pdftitle={Coheteros Progress Report EuRoC 2026},
}

% =========================================================================
% INICIO DEL DOCUMENTO (Debe ir antes de la portada)
% =========================================================================
\begin{document}

% =========================================================================
% PORTADA CENTRADA EN INGLÉS (ESTILO INFORME TÉCNICO OFICIAL)
% =========================================================================
\begin{titlepage}
    % Aseguramos un margen equilibrado para contener todo en una sola página
    \newgeometry{left=2.5cm,right=2.5cm,top=2.5cm,bottom=2.5cm}
    \centering

    % --- LOGOTIPO DE COHETEROS (CENTRADO) ---
    \vspace*{0.5cm}
    \includegraphics[height=2.2cm]{logoSVG_v3_negro.png} 
    
    \vspace{2.5cm}

    % --- CABECERA DE LA COMPETICIÓN ---
    {\large\scshape\color{EtseNavy} European Rocketry Challenge 2026} \\
    \vspace{0.4cm}
    {\large\bfseries PROGRESS REPORT}

    \vspace{1.5cm}

    % --- BLOQUE DE TÍTULO ---
    \vspace{0.6cm}
    {\Huge \bfseries \color{EtseNavy} PROJECT NAOS EVO} \\
    \vspace{0.4cm}
    {\Large \itshape Design, Development, and Flight Readiness Status} \\
    \vspace{1.2cm}

    % --- SUBTÍTULO OBLIGATORIO EUROC ---
    {\large \textbf{Team 07 Progress Report to the 2026 EuRoC}} \\

    \vspace{1cm}

    \begin{figure} [H]
        \centering
        \includegraphics[width=0.5\linewidth]{Cohete_NAOS_Despegue.jpeg}
        \label{fig:placeholder}
    \end{figure}
    
    \vfill % Empuja los metadatos al final de la página de forma equilibrada

    % --- METADATOS INFERIORES (100% INGLÉS) ---
    {\bfseries\large Coheteros} \\
    \vspace{0.15cm}
    {\small Rocketry Association} \\
    {\small School of Engineering} \\
    {\small University of Seville} \\
    \vspace{0.6cm}
    {\small July 9, 2026}

    \vspace*{0.5cm}
    \restoregeometry
\end{titlepage}
\newpage % Fin de la portada, pasamos a la siguiente página limpia

% =========================================================================
% AQUÍ COMIENZA EL CUERPO DEL INFORME (Sin \maketitle)
% =========================================================================

\tableofcontents
\newpage

\listoffigures
\newpage
\listoftables
\newpage

% =========================================================================
% SECTION 1: INTRODUCTION & OBJECTIVES
% =========================================================================
\section{Project Introduction and Objectives}

\subsection{Team Introduction and Project Context}

Coheteros is a student-led experimental rocketry association based at the Escuela Técnica Superior de Ingeniería (ETSI) of the Universidad de Sevilla. The association's primary focus is to advance aerospace engineering education through the practical, hands-on development of launch vehicles. The NAOS EVO represents the association's flagship single-stage launch vehicle, engineered specifically to compete in the 2026 European Rocketry Challenge (EuRoC). The architecture incorporates extensive technological heritage from the previous NAOS launcher iteration.\\

The core technical philosophy driving the team is a strict commitment to 100\% Student Researched and Developed (SRAD) engineering. By minimising the use of Commercial Off-The-Shelf (COTS) components, the team avoids black-box solutions and ensures a fundamental understanding of flight physics. This approach extends deeply into the software domain; rather than relying on standard commercial packages, the team designs, writes, and iterates custom in-house simulators for aerodynamics, propulsion, and Guidance, Navigation, and Control (GNC). \\

Within this context, NAOS EVO serves as an iterative test bench to validate theoretical modelling with empirical flight data. The project challenges students to manage the complete lifecycle of complex subsystems, gathering solid propulsion, active control mechanisms, and parachute manufacturing, transforming academic curiosity into functional, real-world aerospace solutions.

\subsection{Project Goals}

The overarching aim of the NAOS EVO programme is to provide students with an authentic, rigorous aerospace engineering experience. Specific project goals include:
\begin{itemize}
    \item \textbf{Flight Execution:} Successfully launch and recover the NAOS EVO vehicle at EuRoC 2026.
    \item \textbf{Regulatory Compliance:} Navigate complex technical challenges by strictly adhering to real-world regulations, safety standards, and EuRoC design requirements.
    \item \textbf{Empirical Validation:} Iteratively improve in-house simulation tools by feeding them with empirical data gathered during ground tests and actual flight.
    \item \textbf{Professional Transition:} Bridge the gap between academic theory and professional engineering practice by successfully defending custom SRAD systems before expert examiners.
\end{itemize}

\subsection{Mission Objectives}

To guarantee mission success and system validation, the flight profile must satisfy the following quantifiable technical objectives:
\begin{itemize}
    \item \textbf{Apogee Targeting:} Attain an apogee of exactly 3000 metres Above Ground Level (AGL) using the custom APCP solid rocket motor and the active airbrake control system.
    \item \textbf{Recovery Deployment:} Execute a safe, two-stage tender-descender recovery sequence, utilising the clean-gas CO\textsubscript{2} ejection mechanism at the predefined altitudes.
    \item \textbf{Data Acquisition:} Maintain a stable RF telemetry link throughout the flight and log high-speed, non-volatile data for comprehensive post-flight aerodynamic and structural analysis.
\end{itemize}

% =========================================================================
% SECTION 2: CONCEPT OF OPERATIONS (CONOPS)
% =========================================================================
\section{Concept of Operations (CONOPS)}
% Requisito: 1 diagrama de las etapas principales de operación y una breve descripción de texto.

The Concept of Operations for the launch vehicle is governed by the flight controller's state machine, which is segmented into four primary operational phases: Ground, Ascent, Descent, and Abort. This architecture ensures deterministic behaviour and defines failsafe transitions throughout the mission profile.

\begin{figure}[H]
    \centering
    % Placeholder for the provided state machine diagram
    \includegraphics[width=\textwidth]{state_machine_diagram.jpg}
    \caption{Avionics state machine diagram detailing nominal flight phases and abort condition transitions.}
    \label{fig:state_machine}
\end{figure}

As illustrated in Figures \ref{fig:state_machine} and \ref{fig:conops_diagram}, nominal operations progress sequentially based on specific telemetry and timing triggers. The transition logic and state actions are detailed in Table \ref{tab:conops_states}.

\begin{figure}[H]
    \centering
    % Placeholder for the provided state machine diagram
    \includegraphics[width=\textwidth]{Pictures/CONOPS_Diagram.png}
    \caption{NAOS EVO launcher Concept of Operation Diagram}
    \label{fig:conops_diagram}
\end{figure}

\begin{table}[H]
    \centering
    \caption{Flight controller state machine transition logic and nominal phase behaviours.}
    \label{tab:conops_states}
    \begin{tabularx}{\textwidth}{|l|l|X|}
        \hline
        \rowcolor{EtseNavy}
        \textbf{\textcolor{white}{Phase}} & \textbf{\textcolor{white}{State}} & \textbf{\textcolor{white}{Condition \& Behaviour}} \\ \hline
        Ground & IDLE & Pseudo-state initiated upon hardware switch ON. \\ \hline
        \rowcolor{LightGrayTable}
        Ground & INIT & Avionics initialisation sequence. Transitions to CALIBRATION upon receiving the remote GO signal via telemetry. \\ \hline
        Ground & CALIBRATION & Sensor baseline establishment. Transitions to PRELAUNCH once sensors are confirmed calibrated. \\ \hline
        \rowcolor{LightGrayTable}
        Ground & PRELAUNCH & SD logging activation. Awaits IGNITION detection algorithm trigger. \\ \hline
        Ascent & BOOST & Active thrust phase. \\ \hline
        \rowcolor{LightGrayTable}
        Ascent & COASTING & Burn-out. Active control disengaged. Triggered by accelerometer-based motor burn-out detection. \\ \hline
        Ascent & ACTIVE CONTROL & Active control algorithms engaged after passing a predefined altitude constraint. \\ \hline
        \rowcolor{LightGrayTable}
        Descent & APOGEE & Drogue parachute deployment, triggered by barometric and inertial apogee detect  ion. \\ \hline
        Descent & MAIN PARACHUTE & Main parachute deployment, triggered by a secondary descent altitude constraint. \\ \hline
        \rowcolor{LightGrayTable}
        Descent & LANDED & Touchdown detection. Disables non-essential sensors and high-current draw systems to preserve battery life for recovery telemetry. \\ \hline
    \end{tabularx}
\end{table}

Anomalies detected in telemetry or control loops force an immediate transition to the respective Abort state. A \textbf{Ground Abort} halts the launch sequence. A \textbf{Descent Abort} overrides nominal altitude logic to force immediate parachute deployment, ensuring the vehicle returns within range limits at a safe velocity.
\textbf{Ascent Abort} states are omitted as the rocket will continue its ballistic path and transition to the respective descent states.


% =========================================================================
% SECTION 3: SYSTEM CONCEPT AND DESIGN
% =========================================================================
\section{System Concept and Design}



\subsection{General Arrangement}
% Requisito: Diagrama/plano general del cohete y 1 o 2 párrafos explicativos.
\begin{wrapfigure}[32]{r}{0.43\textwidth} % <-- El [40] fuerza la altura exacta en líneas
    \centering
    \vspace{-10pt} % Ajuste fino opcional para subirla un poco si hace falta
    \includegraphics[width=1\linewidth]{Pictures/System_Arrangement_Diagram.png}
    \caption{General arrangement of the NAOS EVO launch vehicle, detailing main internal subsystems and modular sections.}
    \label{fig:NAOS_EVO_General_Arrangement}
\end{wrapfigure}

NAOS EVO is a single-stage, solid-propellant launch vehicle engineered with a highly modular architecture. Marking a significant structural milestone for the association, it is the first vehicle developed by Coheteros to utilise carbon-fibre composites as the primary structural and load-bearing material. This material selection ensures high structural stiffness while strictly adhering to the rigorous integration standards expected in high-power rocketry.\\

Propulsion is provided by a custom Student Researched and Developed (SRAD) solid motor, operating on an Ammonium Perchlorate Composite Propellant (APCP) formulation with an HTPB binder. For the descent phase, the vehicle features a tender-descender dual-deployment recovery system. This architecture consists of a drogue and a main parachute. Upon reaching a safe predefined altitude, the vehicle's mass acting against the aerodynamic drag of the drogue parachute provides the mechanical extraction force required to safely deploy the main canopy.\\

Additionally, the vehicle incorporates an SRAD active control system. This airbrake mechanism modulates aerodynamic drag during the coast phase, fulfilling a dual purpose: ensuring precise 3000-metre apogee targeting and providing a reliable risk mitigation method to scrub kinetic energy during emergency abort scenarios.

\vspace{1em}
\newpage
% =========================================================================
% SUBSECTION 3.2: Dimensions and Mass Estimates
% =========================================================================
\subsection{Dimensions and Mass Estimates}
% Requisito: Plano acotado y/o tabla detallada con estimaciones de peso.

The geometric parameters and mass distribution of the NAOS EVO vehicle are strictly controlled to ensure aerodynamic stability and adherence to the target apogee. Figure \ref{fig:dimensioned_drawing} provides the overall dimensioned layout of the vehicle, highlighting the relative positions of the primary subsystems.

\begin{figure}[H]
    \centering
    % Placeholder for the dimensioned drawing
    \includegraphics[width=0.8\textwidth]{Figure 3.png}
    \caption{Dimensioned drawing of NAOS EVO, highlighting the overall layout, Centre of Gravity (CoG), and Centre of Pressure (CP) from RocketPy.}
    \label{fig:dimensioned_drawing}
\end{figure}

Aerodynamically, the launch vehicle features an outer diameter of 91 mm. Based on the current layout, the Centre of Pressure (CP) is located at 1711.8 mm from the nose tip, providing a nominal static margin of 3.24 calibres off the launch rail. Table \ref{tab:mass_layout_breakdown} details the preliminary mass budget and longitudinal layout broken down by functional subsystem, providing a comprehensive overview of the vehicle's structural distribution and CoG.


\begin{table}[H]
    \centering
    \caption{Preliminary mass and longitudinal layout breakdown by subsystem.}
    \label{tab:mass_layout_breakdown}
    \renewcommand{\arraystretch}{1.3}
    \begin{tabularx}{\textwidth}{|>{\raggedright\arraybackslash}X|
                                >{\centering\arraybackslash}p{2.8cm}|
                                >{\centering\arraybackslash}p{2.3cm}|
                                >{\centering\arraybackslash}p{3.0cm}|}
        \hline
        \rowcolor{EtseNavy}
        \textbf{\textcolor{white}{Subsystem}} & \textbf{\textcolor{white}{Fuselage Length [mm]}} & \textbf{\textcolor{white}{Mass [kg]}} & \textbf{\textcolor{white}{CoG Position\textsuperscript{1} [mm]}} \\
        \hline
        Nosecone\textsuperscript{2} & 273 & 0.142 & 261.1 \\
        \hline
        \rowcolor{LightGrayTable}
        Recovery System (excluding nosecone) & 575 & 1.007 & 532.4 \\
        \hline
        Payload & 270 & 1.077 & 982.7 \\
        \hline
        
        % --- INICIO BLOQUE AGRUPADO AVIONICS/CONTROL ---
        % Ambas filas van en el mismo color para que la celda combinada quede limpia
        \rowcolor{LightGrayTable}
        Avionics & & 0.704 & 1255.8 \\
        \cline{1-1}\cline{3-4}
        \rowcolor{LightGrayTable}
        Active Control System & \multirow{-2}{=}{\centering 470} & 0.649 & 1496.3 \\
        \hline
        % --- FIN BLOQUE AGRUPADO ---
        
        Motor (Loaded) & 605 & 4.119 & 1866.8 \\
        \hline
        \rowcolor{LightGrayTable}
        Fins & --- & 0.304 & 2144.3 \\
        \hline
        Other\textsuperscript{3} & --- & 2.858 & --- \\
        \noalign{\hrule height 1.2pt}
        \rowcolor{EtseNavy}
        \textbf{\textcolor{white}{Total Vehicle}} & \textbf{\textcolor{white}{2192}} & \textbf{\textcolor{white}{10.86}} & \textbf{\textcolor{white}{1417.5}} \\
        \hline
    \end{tabularx}
    
    \vspace{2mm}
    \footnotesize{
    \begin{flushleft}
    \textsuperscript{1} Measured from the nose tip. Values are preliminary CAD estimates and subject to refinement.\\
    \textsuperscript{2} Length corresponds to the visible external profile after assembly, excluding internal coupling shoulders.\\
    \textsuperscript{3} Includes fuselage sections, section-joining bulkheads, and minor structural integration hardware.
    \end{flushleft}
    }
\end{table}



\subsection{Main Performance Figures}
% Requisito: Tabla resumen de las métricas de rendimiento clave.

The vehicle's theoretical flight performance and internal motor ballistics are summarised in Table \ref{tab:performance_figures}. Following initial static fire tests that successfully qualified the propellant, the motor's grain configuration was iterated to a 6-Bates geometry (82 mm length, 20 mm core) combined with a 14 mm nozzle throat. This design ensures the vehicle meets the strict 3000-metre apogee target.

\begin{table}[H]
    \centering
    \caption{Key flight performance and internal ballistics parameters.}
    \label{tab:performance_figures}
    \begin{tabularx}{\textwidth}{|X|X|}
        \hline
        \rowcolor{EtseNavy}
        \textbf{\textcolor{white}{Parameter}} & \textbf{\textcolor{white}{Estimated Value / Target}} \\ \hline
        Target Apogee & 3000 m AGL \\ \hline
        \rowcolor{LightGrayTable}
        Maximum Velocity & Mach 0.9 \\ \hline 
        Propellant Mass ($m_{prop}$) & 1.689 kg \\ \hline
        \rowcolor{LightGrayTable}
        Total Impulse ($I_t$) & 3521 Ns \\ \hline
        Maximum Thrust ($T_{max}$) & 1287 N \\ \hline
        \rowcolor{LightGrayTable}
        Average Thrust ($T_{mean}$) & 1136 N \\ \hline
        Burn Duration ($t_b$) & 3.05 s \\ \hline
        \rowcolor{LightGrayTable}
        Specific Impulse ($I_{sp}$) & 212 s \\ \hline
        Maximum Chamber Pressure ($P_0$) & 59 bar \\ \hline
        \rowcolor{LightGrayTable}
        Main Parachute Deployment Altitude & 450 m AGL \\ \hline 
    \end{tabularx}
\end{table}


\subsection{Main Systems Description}
% Requisito: Breve introducción de transición antes de detallar cada subsistema.

This section details the individual engineering designs, mechanical architectures, and operational rationales of the primary subsystems integrated within the NAOS EVO launch vehicle.

\subsubsection{Avionics Subsystem}\label{ssubsec:avionics_design}
The avionics architecture employs a modular, distributed topology partitioned across three custom-designed printed circuit boards (PCBs) to minimise radio frequency (RF) noise coupling and simplify physical integration within the electronics bay. The primary MCU \& Sensor Board serves as the central processing unit of the launch vehicle, utilising an STM32 microcontroller to poll flight dynamics, log high-speed data to non-volatile storage, and execute real-time guidance and control algorithms for the active airbrake system. \\

The core flight controller firmware is written in C using the STM32CubeMX hardware abstraction layer within STM32CubeIDE. System execution is governed by FreeRTOS implementing a deterministic finite state machine. This real-time operating system structure manages ground station command processing, down-link telemetry transmission, multi-sensor data acquisition, Kalman filter state estimation, and conditional transitions between discrete flight phases.

\begin{figure}[H]
    \centering
    \begin{minipage}{0.45\textwidth}
        \centering
        \includegraphics[height=5.5cm, keepaspectratio]{avionica 1.png}
        \\[0.15cm] \small (a) Primary/Sensor Face
    \end{minipage}
    \hfill
    \begin{minipage}{0.45\textwidth}
        \centering
        \includegraphics[height=5.5cm, keepaspectratio]{avionica 2.png}
        \\[0.15cm] \small (b) Communications Face
    \end{minipage}
    \vspace{0.3cm}
    \caption{Three-dimensional CAD assembly of the internal avionics bay stack, highlighting the longitudinal structural threaded rods, PCB mounting brackets, and the integrated upper cylindrical bulkhead mechanism. This avionics bay configuration may change.}
    \label{fig:avionics_cad_assembly}
\end{figure}

RF telemetry and GNSS tracking are physically isolated on a dedicated Communications Board operating on an 868\,MHz LoRa link to prevent electromagnetic interference with sensitive analogue sensor traces. Power regulation is managed by an independent Power Distribution Board (PDB), incorporating high-efficiency switching DC-DC converters to supply isolated 5\,V and 3.3\,V rails.\\

Ground operations and safety compliance are enforced through two distinct physical interfaces: a primary mechanical switch linking the main battery to the PDB, and an independent, isolated arming switch for the pyrotechnic deployment charges. This dual-switch configuration ensures complete galvanic isolation of the recovery mechanism prior to formal launch clearance. \\

To satisfy EuRoC safety and reliability requirements, a fully redundant flight computer operates concurrently in a hot-standby configuration. This secondary system utilises an independent STM32 MCU and an identical, isolated sensor suite comprising an inertial measurement unit (IMU), a barometer and a magnetometer. \\

A hardware power multiplexer provides fault tolerance for the 5\,V airbrake servo rail, while a solid-state electrical switch handles PWM signal routing. The secondary computer executes the state machine and Kalman filter algorithms in parallel with the primary unit. This setup enables immediate control arbitration of the airbrake system if a primary hardware fault is detected, while continuously logging baseline telemetry to an independent SD card for post-flight analysis.

\begin{figure}[H]
    \centering
    \includegraphics[height=4.2cm, keepaspectratio]{BSB.jpg}
    \caption{Manufactured Central MCU \& Sensor Board (CHT\_BSB\_v1.0) displaying the central STM32 microcontroller footprint, integrated microSD slot for high-speed non-volatile data logging, debug interfaces, and local tactile switches.}
    \label{fig:pcb_baseboard}
\end{figure}

\begin{figure}[H]
    \centering
    \includegraphics[height=4.2cm, keepaspectratio]{PWB.jpg}
    \caption{Manufactured Power Distribution Board (PDB) showing the integrated high-current input connector, status LED indicators, and three discrete Omron 5\,V DC relays dedicated to drogue, main parachute deployment actuation, and servo rail isolation.}
    \label{fig:pcb_powerboard}
\end{figure}

\subsubsection{Active Airbrake System}
\label{ssubsec:airbrakes}

The rocket features an active airbrake system for precise apogee targeting through drag modulation. The mechanism, almost entirely machined in 7075-T6 aluminium, consists of four flat petals that deploy radially via a servo-actuated slider-crank mechanism. Additionally, all custom aluminium parts are designed to be the same thickness for cost reduction and ease of manufacture, and the mechanism features linear rail guides to reduce friction and prevent jamming.\\

Airbrake actuation is governed by a closed-loop guidance and control algorithm that computes reference drag coefficients in real time and commands the deployment level required to achieve them. The closed-loop guidance algorithm works by periodically computing predicted apogee height from current vehicle state as a function of drag coefficient, and root finding is then used to obtain a reference drag coefficient. Using pre-computed lookup tables, the job of the control algorithm is to then deploy the airbrakes to reach the reference drag coefficient.

\begin{figure}[H]
    \centering
    \includegraphics[width=\textwidth]{GNC_loop_diagram.png}
    \caption{Conceptual block diagram showing the full Guidance, Navigation and Control (GNC) loop of the vehicle.}
    \label{fig:GNC_diagram}
\end{figure}

Since deployment is intended at 1500 metres AGL in a nominal flight, which is some time after burnout, when the rocket will already have lost some speed, the petals are positioned at two different axial locations along the airframe, thereby maximizing the effective available drag area for greater control authority.\\

CFD and finite-element analyses have been performed with the expected flight envelope in mind and the results show stresses to be well within the structural limits of the materials used. In addition, the linear guides were tested slightly above expected flight loads and did not experience considerably higher friction or jamming. The results show that the servo motor can provide more than four times the torque required to actuate the airbrakes at maximum dynamic pressure, which is the moment when friction would be at its highest.\\

During the mechanical design phase, the main challenges encountered were size and budget constraints. For that reason, the team chose to rely on metal machining as little as possible and instead bought off-the-shelf components whenever feasible. Regarding size constraints, due to the relatively small diameter of the rocket, the final design is the result of several iterations and adjustments. The peculiar asymmetric location of the linear rail guides is a direct consequence of that, given that a symmetric configuration would make the rails impact the central shaft. In order to counteract the possible roll moment contributions that could stem from the rail and wake interactions due to that asymmetry, the top and bottom axial locations of the airbrakes have their rails located at opposing locations, which has been found to mitigate this effect.

\begin{figure}[H]
    \centering
    \includegraphics[width=0.5\textwidth]{Airbrakes_top_view.jpeg}
    \caption{Top view of the airbrake mechanism.}
    \label{fig:Airbrake_top}
\end{figure}

\begin{figure}[H]
    \centering
    \begin{minipage}{0.45\textwidth}
        \centering
        \includegraphics[height=9cm, keepaspectratio]{Airbrakes_stowed.jpeg}
        \\[0.15cm] \small (a) Stowed airbrakes
    \end{minipage}
    \hfill
    \begin{minipage}{0.45\textwidth}
        \centering
        \includegraphics[height=9cm, keepaspectratio]{Airbrakes_deployed.jpeg}
        \\[0.15cm] \small (b) Deployed airbrakes
    \end{minipage}
    \vspace{0.3cm}
    \caption{CAD assembly of the airbrake mechanism as currently being manufactured.}
    \label{fig:airbrake_views}
\end{figure}

As for the electronics, the servo motor is interfaced through a hardware multiplexer to two independent flight computers to ensure command redundancy and is powered by two independent five-volt supply lines through a power multiplexer circuit. That way, if either battery dies or either flight computer stops responding, there will be an additional working computer that can resume control.\\

The flight computers are set up in an active-standby configuration, where the standby (backup) computer constantly monitors a heartbeat from the active (primary) computer. If the primary computer stops responding, the backup computer commands the hardware multiplexer to hand control over to itself, whereas if the backup computer stops working, no command is sent to the hardware multiplexer and the primary computer remains in control.\\

Lastly, the flight computers constantly poll the vehicle state estimation to determine if it is flying on a nominal trajectory. In the case that it is not, a full deployment of the airbrakes is commanded to slow the vehicle down as fast as possible, thus mitigating safety risks.

\subsubsection{Propulsion System}

\paragraph{Propellant}
The primary performance requirement for the NAOS EVO launch vehicle is to achieve a target apogee of at least $3000\,\text{m}$. Two solid propellant candidates were evaluated to satisfy the non-toxicity criteria imposed by the competition rules: Sorbitol-Potassium Nitrate (KNSB) and Ammonium Perchlorate Composite Propellant (APCP). While both choices are chemically viable, APCP was selected based on the following engineering trade-offs:

\begin{itemize}
    \item \textbf{Specific Impulse ($I_{sp}$):} APCP provides a significantly higher specific impulse compared to KNSB. Flight simulations with an identical propellant mass envelope demonstrate that the APCP variant yields nearly twice the apogee altitude of its KNSB counterpart.
    \item \textbf{Mechanical Properties:} KNSB grains are inherently brittle and highly susceptible to structural cracking during handling or ignition pressure spikes, which can lead to catastrophic chamber overpressurisation. Conversely, APCP utilises a Hydroxyl-terminated polybutadiene (HTPB) binder matrix, ensuring excellent structural flexibility. This selection eliminates the chemical caramelisation risks and hygroscopic degradation associated with sugar-based propellants.
\end{itemize}

The final propellant formulation consists of ammonium perchlorate oxidiser, aluminium powder fuel, and an HTPB binder system, supplemented with burn-rate catalysts and curing agents. KNSB remains designated as the secondary backup alternative.

\begin{figure}[H]
    \centering
    \includegraphics[width=0.35\textwidth]{Pictures/propellant_grain.png}
    \caption{Manufactured APCP solid propellant grain segment with a cylindrical core geometry.}
    \label{fig:propellant_grain}
\end{figure}

\paragraph{Combustion Chamber}
The structural casing must withstand transient internal pressures and severe thermal fluxes without leakage while maintaining an acceptable safety margin. Despite its adverse impact on the mass budget due to a higher density, stainless steel was selected over aluminium for the motor casing to maximise thermal-structural margins. Standard stainless steel tubing also offers excellent commercial availability and cost-efficiency.\\

Based on an estimated APCP propellant mass of $1.6\,\text{kg}$ to $1.8\,\text{kg}$ configured in a BATES grain geometry, outer diameters between $60\,\text{mm}$ and $80\,\text{mm}$ were simulated. The $60\,\text{mm}$ diameter envelope was selected to optimise the rocket's overall aerodynamic cross-section. Following preliminary structural characterisation, the casing geometry was finalised with an inner diameter of approximately $60\,\text{mm}$ and an overall length of $525\,\text{mm}$. 

\begin{figure}[H]
    \centering
    \includegraphics[width=0.9\textwidth]{Pictures/Plano 0 - Motor_as built.jpg}
    \caption{Preliminary technical cross-sectional and isometric blueprint of the stainless steel combustion chamber assembly.}
    \label{fig:chamber_blueprint}
\end{figure}

\paragraph{Nozzle and Aft Closure}
The critical dimension governing nozzle performance is the throat diameter. Thermodynamic simulations matching the chamber volume and propellant mass indicate that a throat diameter between $12.5\,\text{mm}$ and $14\,\text{mm}$ is required to maintain a steady-state chamber pressure ($P_0$) between $40\,\text{bar}$ and $60\,\text{bar}$. \\

Based on empirical internal ballistics data, the initial throat diameter is set at $14\,\text{mm}$, subject to minor iterative refinements. The divergent section utilises a truncated cone profile to minimise machining complexity and production costs. The exit diameter is dimensioned to ensure the nozzle operates close to an adapted state for the maximum duration of the motor burn phase. The aft closure is machined from stainless steel to ensure material homogeneity with the combustion chamber walls.\\

\paragraph{Ignition}
Motor ignition is initiated via a pyrotechnic igniter paired with a commercial electric match (e-match). The exact pyrotechnic chemical composition is selected based on subscale qualification tests. This subsystem remains subject to ongoing optimisation to ensure reliable, instantaneous ignition sequences on the launch rail.


\subsubsection{Recovery System}

The recovery system is based on a two-stage parachute configuration, with the drogue parachute deployed near apogee and the main parachute deployed at approximately 450\,m Above Ground Level (AGL), in accordance with mission requirements.\\

The drogue parachute is hemispherical, with a nominal diameter of 0.77\,m, and is designed to withstand an estimated deployment speed of 34\,m/s.\\

The main parachute utilises a ringsail design, replacing the cross-type parachutes employed in previous campaigns. This configuration was selected because a ringsail canopy provides superior aerodynamic stability, reduces oscillation due to its vented ring structure, produces lower opening shock loads, and performs better during high-speed deployment.\\ 

Due to the increased manufacturing complexity associated with the ringsail geometry, fabrication is outsourced to a specialised aerospace sponsor to ensure strict compliance with precise aerodynamic and structural tolerances. The main parachute is designed for an opening speed of 34\,m/s (with margins) to ensure a safe terminal descent rate of 9\,m/s, featuring a nominal diameter of 1.85\,m.\\

% --- INICIO DEL WRAPFIGURE (Ojiva vertical a la derecha) ---
\begin{wrapfigure}[32]{r}{0.4\textwidth}
    \centering
    \vspace{-0.3cm} % Alínea la parte superior de la imagen con el inicio del párrafo
    \includegraphics[width=1\linewidth]{Pictures/Recovery_System_1.png}
    \caption{Nosecone profile section detailing the integrated upper recovery compartment.}
    \label{fig:Recovery_CAD_1-1}
    \vspace{-0.6cm} % Evita que empuje el texto de abajo
\end{wrapfigure}

The final selection of parachute geometries and dimensions will be completed after experimental validation of the aerodynamic data used in simulations. If a change in parachute configuration becomes necessary during final qualification, a standard cross-type canopy remains designated as the backup option.\\

The separation event at apogee is achieved by detaching the ogive from the main body, thereby initiating drogue parachute deployment. During ascent, the ogive remains mechanically secured to the airframe, enclosing the drogue compartment.\\

The first deployment stage relies on a clean-gas CO\textsubscript{2} pressurisation mechanism. Upon receiving an electrical signal from one of the flight computers, a small pyrotechnic charge actuates a puncture pin, driving it against a sealed CO\textsubscript{2} cartridge. The released gas increases the internal pressure within the drogue compartment, forcing separation of the ogive from the main body and enabling drogue parachute extraction.\\

After deployment, the drogue parachute remains connected to the main body via a cord incorporating a tender-descender mechanism. Upon reaching the predefined main deployment altitude (approximately 450\,m AGL), a black powder charge activates the tender-descender, severing the primary cord. This action transfers the mechanical load to a secondary connection linked to the recovery bulkhead. Once tensioned, the bulkhead is extracted, which in turn deploys the main parachute, as both elements are mechanically linked. This staged configuration ensures controlled extraction and minimises the risk of line entanglement or premature main canopy deployment.

% --- PLANO HORIZONTAL FLOTANTE (Colocado al final para evitar colisiones) ---
\begin{figure}[H]
    \centering
    \vspace{0.3cm}
    \includegraphics[width=0.95\textwidth]{Pictures/Recovery System_2.jpeg} 
    \caption{CAD layout of the internal recovery system packaging and structural interfaces.}
    \label{fig:Recovery_CAD_1}
\end{figure}

\subsubsection{Airframe}
To evaluate the aerodynamic performance and stability of the launch vehicle, a custom low-fidelity aerodynamic prediction software named StarCom ({Figure}~\ref{fig:starcom_logo}) was utilized. Currently under active development within the association, StarCom integrates a proprietary graphical layer built in MATLAB over the robust analytical core engine of Digital DATCOM, providing a reliable computational foundation to rapidly iterate, analyze, and validate complex aerodynamic configurations.

\begin{figure}[H]
    \centering
    \includegraphics[width=0.4\textwidth]{Pictures/STARCOM.jpg}
    \caption{Official software insignia for the STARCOM aerodynamic prediction suite.}
    \label{fig:starcom_logo}
\end{figure}

The launch vehicle is engineered with a highly slender airframe, characterized by a length-to-diameter ($L/D$) fineness ratio of 24. To manage the specific aerodynamic characteristics inherent to high fineness ratios, a nominal design static margin of 3.5 calibres was established (Figure~\ref{fig:static_margin_plots}). This margin corresponds to 12\% of the vehicle's total longitudinal length, a value well-supported by experimental aerospace literature to prevent both under-stability and excessive over-stability. As depicted in the computational simulations, this configuration guarantees positive restoring moments and a predictable flight trajectory, maintaining robust stability margins for angles of attack (AoA) up to 20$^\circ$ across the operational velocity regime.

\begin{figure}[H]
    \centering
    \includegraphics[width=1\textwidth]{Pictures/SM_Plot.jpeg}
    \caption{Static margin (SM) variation profiles as a function of Mach number and angles of attack (AoA).}
    \label{fig:static_margin_plots}
\end{figure}

Regarding the vehicle's external lifting and aerodynamic surfaces (Figure~\ref{fig:datcom_3d_model}), a Von Kármán nose cone profile was selected to minimize wave drag coefficient within the transonic flight regime. The stabilizing trapezoidal fins \cite{fleeman_tactical_2006} implement a flat-plate airfoil geometry. Numerical analysis confirms that within the designated flight envelope, the flat plate retains optimal aerodynamic properties by effectively acting as a sharp-edged supersonic airfoil. This design choice satisfies all stability requirements while significantly streamlining structural manufacturability and reducing assembly tolerances.

\begin{figure}[H]
    \centering
    \includegraphics[width=0.8\textwidth]{Pictures/DATCOM_vehicle_3Dplot.jpeg}
    \caption{Three-dimensional aerodynamic surface plot and geometric discretization generated for digital DATCOM parsing.}
    \label{fig:datcom_3d_model}
\end{figure}



% =========================================================================
% SECTION 4: SPECIAL DESIGN FEATURES AND MATERIALS
% =========================================================================
\section{Special Design and Manufacturing Features}

\subsection{Detailed Special Design Features}
% Requisito: Diagramas/esquemas y 2 o 3 párrafos de texto.

NAOS EVO uses custom navigation systems to estimate its state in flight. Its flight computers run an Error-State Kalman Filter (ESKF) developed by the team in order to track 10 states: 3D position, 3D velocity and attitude quaternion. The ESKF fuses noisy data from several sources, which in this case are an accelerometer, a gyroscope, a magnetometer and a barometer.\\

An especially interesting feature of the ESKF is that sensors which are not traditionally used, for example, for attitude estimation, like the barometer, can still provide some attitude information, given that attitude can be slightly observable through barometric data due to cross-covariance relationships.

\subsection{Materials and Methods of Manufacture}
% Requisito: 1 paragraph or a list

The material selection for NAOS EVO is primarily driven by accessibility and manufacturing feasibility. The vehicle relies on widely available, cost-effective materials, combining in-house fabrication techniques, such as FDM 3D printing for internal components and wet lamination for composites, with standard external machining processes for metal parts, such as the stainless steel aft closure. Table \ref{tab:materials_processes} outlines the primary materials employed and their corresponding methods of manufacture.
\begin{table}[H]
    \centering
    \caption{Primary materials and associated manufacturing processes.}
    \label{tab:materials_processes}
    \begin{tabularx}{\textwidth}{|>{\centering\arraybackslash}X|>{\centering\arraybackslash}X|}
        \hline
        \rowcolor{EtseNavy} 
        \textbf{\textcolor{white}{Material}} & \textbf{\textcolor{white}{Manufacturing Process}} \\ \hline
        Acrylonitrile Butadiene Styrene (ABS) & FDM 3D Printing \\ \hline
        \rowcolor{LightGrayTable} 
        Polylactic Acid (PLA) & FDM 3D Printing \\ \hline
        Aluminium Alloys & CNC Machining \\ \hline
        \rowcolor{LightGrayTable} 
        AISI 304 Stainless Steel & Lathe Machining (Turning) \\ \hline
        Glass-fibre Composite & Wet Lamination \\ \hline
        \rowcolor{LightGrayTable} 
        Carbon-fibre Composite & Wet Lamination \\ \hline
        Electronic Components, PCBs & Reflow Soldering, hand soldering \\ \hline
    \end{tabularx}
\end{table}

\subsection{Differentiating and Unique Characteristics}
% Requisito: 1 o 2 párrafos más esquema explicativo.

The NAOS EVO launch vehicle incorporates several proprietary innovations that distinguish it from conventional COTS-reliant systems, particularly in guidance, navigation, and control. Trajectory modulation is achieved via a custom active airbrake mechanism, which utilises a central rotating shaft to deploy aluminium petals radially and simultaneously. This mechanical synchronisation ensures symmetric drag distribution during the coast phase, governed by a Model Predictive Control (MPC) inspired algorithm. Supporting this architecture, the avionics suite runs an Error-State Kalman Filter (ESKF). By fusing data from the IMU, barometer, and magnetometer, the system estimates not only position but also vehicle attitude in real-time, establishing the foundational architecture for active attitude control in future vehicles.\\

In addition to aerodynamic control, the vehicle features highly integrated, student-researched propulsion and recovery mechanisms. The propulsion module is driven by a proprietary Ammonium Perchlorate Composite Propellant (APCP) solid motor. This vertical integration encompasses a custom chemical formulation, manufacturing procedure, and static fire testing protocol, allowing for precise characterisation of the thrust curve and optimised grain geometries. For recovery, the vehicle departs from traditional pyrotechnic ejection charges in favour of a clean-gas CO\textsubscript{2} system. Avoiding explosive combustion within the confined recovery bay mitigates structural stress peaks, eliminates the risk of thermal damage to the parachutes, and significantly enhances operational safety during ground handling and launch preparation.

\begin{figure}[H]
    \centering
    % Placeholder for a schematic/drawing of one of the unique features (e.g., the Airbrake CAD or CO2 mechanism)
    \includegraphics[width=1\textwidth]{Pictures/Diagram_NAOS_EVO_speacial_features.png}
    \caption{Differentiating and unique characteristics Drawing}
    \label{fig:differentiating_features}
\end{figure}

% =========================================================================
% SECTION 5: CRITICALITIES AND TESTING STATUS
% =========================================================================
\section{Criticalities, Risk Analysis, and Experimental Validation}


\subsection{Expected Difficulties and Criticalities}
% Requisito: 3 o 4 párrafos identificando los puntos más críticos.
To start with, the clean gas recovery ejection will be implemented for the first time, and as such will require thorough testing and integration oversight.\\

Another issue to assess is the extraction of the main parachute, which depends entirely on the successful inflation of the drogue. A weak drogue deployment, due to low dynamic pressure or entanglement, would prevent the main parachute from being extracted from the bulkhead.\\

Ongoing work focuses on ground validation of the CO$_2$ pressurization system, structural verification of bulkhead release loads, and reliability testing of the tender-descender mechanism. Special emphasis is placed on repeatability, integration efficiency, and minimizing turnaround time during launch operations. Flight testing with an instrumented test vehicle is planned to validate descent rates and opening shock loads under representative conditions.\\

Additionally, the airbrake system is also being used for the first time. Therefore, it will be rigorously load tested to ensure symmetric petal deployment and a low probability of jamming at maximum expected flight loads. Finally, HIL (Hardware-In-The-Loop) simulations will be performed to test the system’s performance in a flight-like environment. 


% =========================================================================
% SUBSECTION 5.2: Tests Performed and Outcomes
% =========================================================================
\subsection{Tests Performed and Outcomes}
\label{Tests_Performed} 
% Requisito: 3 o 4 párrafos que listen los ensayos realizados y sus resultados.

\textbf{Propulsion and Static Fire Testing:} Ground static hot fire tests of the custom solid rocket motor have been successfully conducted to qualify the Ammonium Perchlorate Composite Propellant (APCP) formulation and validate the structural integrity of the stainless-steel casing (Figure \ref{fig:static_fire_test}). These campaigns verified propellant combustion characteristics under transient internal pressures, confirming reliable ignition and stable burning rates. The empirical pressure ($P_0$) and thrust curves recorded across multiple test dates (Figure \ref{fig:thrust_pressure_curves}) demonstrated high repeatability. Although initial subscale tests did not meet the target 3000-metre apogee due to non-representative test parameters, the gathered ballistic data directly informed a design iteration. This led to the definitive 6-Bates grain configuration (82\,mm length, 20\,mm core) with a 14\,mm nozzle throat, yielding the optimised theoretical performance validated in the subsequent test campaign on 28/06/2026 (Figure \ref{fig:thrust_pressure_curve_28-06-2026}).

\begin{figure}[H]
    \centering
    \includegraphics[width=0.32\textwidth]{Pictures/static_fire.png}
    \caption{Ground static hot fire testing of the SRAD APCP motor unit.}
    \label{fig:static_fire_test}
\end{figure}

\begin{figure}[H]
    \centering
    \includegraphics[width=1\textwidth]{Pictures/thrust_pressure_curves.png}
    \caption{Empirical thrust (kg) and chamber pressure (bar/atm) curves versus time (s) recorded during static fire testing.}
    \label{fig:thrust_pressure_curves}
\end{figure}

\begin{figure}[H]
    \centering
    \includegraphics[width=1\textwidth]{Pictures/thrust_pressure_curve_28-06-2026.png}
    \caption{Theoretical versus empirical thrust and pressure comparison plots for the iterated 28/06/2026 test.}
    \label{fig:thrust_pressure_curve_28-06-2026}
\end{figure}

\newline

\textbf{Recovery System Ejection and Structural Validation:} The clean-gas CO\textsubscript{2} ejection mechanism has undergone ground deployment testing, capturing successful separation dynamics (Figure \ref{fig:recovery_tests_combined}a). These tests utilised additively manufactured PLA and ABS hardware, 12\,g CO\textsubscript{2} cartridges, and 0.5\,g black powder charges to validate the separation kinetics of the ogive nosecone (Figure \ref{fig:Recovery_CAD_1}). Ground deployment tests verified that the custom pressurisation mechanism successfully generates extraction forces exceeding five times the nominal drogue parachute weight, satisfying standard aerospace recovery design criteria \cite{knacke1992}. Structural load testing also successfully validated the tender-descender mechanism under simulated flight loads. In-house parachute manufacturing quality has been verified through prototype fabrication of the hemispherical canopy (Figure \ref{fig:recovery_tests_combined}b), with aerodynamic characterisation scheduled for wind-tunnel testing to quantify opening times, dynamics, and drag coefficients ($S\cdot C_{D_0}$). The second iteration of the deployment mechanism improved upon the first by increasing the shaft-hole interference at the retention interface. This modification raised initial friction, delaying separation until the CO\textsubscript{2} cartridge had fully discharged, thereby producing a more energetic and repeatable release. Post-test analysis of system geometry and separation loads is currently driving the refinement of flight-ready materials and tighter interface tolerances.

\begin{figure}[H]
    \centering
    % Usamos [b] para alinear las bases y que los subtítulos queden en línea
    \begin{minipage}[b]{0.54\textwidth}
        \centering
        \includegraphics[height=5.0cm, keepaspectratio]{Pictures/Recovery_System_test.jpg}
        \\[0.15cm] \small (a) Recovery system deployment test
    \end{minipage}
    \hspace{0.6cm} % Reemplazamos \hfill por un espacio fijo para centrar ópticamente el bloque
    \begin{minipage}[b]{0.36\textwidth}
        \centering
        \includegraphics[height=5.0cm, keepaspectratio]{Pictures/parachute.jpg}
        \\[0.15cm] \small (b) Hemispherical parachute manufacturing
    \end{minipage}
    \vspace{0.3cm}
    \caption{Experimental validation of the clean-gas ejection mechanism and the drogue parachute manufacturing process.}
    \label{fig:recovery_tests_combined}
\end{figure}

\textbf{Avionics Hardware-In-The-Loop (HIL) and Power Integration:} To validate the deterministic behaviour of the flight controller, HIL testing was executed by injecting simulated flight telemetry via a serial UART interface, bypassing the physical sensors. Concurrently, comprehensive functional testing of the Power Distribution Board (PDB) and MCU Baseboard verified power routing and signal integrity. The logic constraints correctly isolated and energised the 5\,V Omron DC relays for parachute pyrotechnic simulation lines without inducing voltage transients. Additionally, the primary MCU successfully routed stable Pulse Width Modulation (PWM) signals to the active airbrake servos, demonstrating complete electrical isolation until explicit command authorisation was granted. The 868\,MHz LoRa transceiver was also verified, maintaining nominal link budgets and packet transmission stability.
\newline

\textbf{Software, GNC, and Data Validation:} Guidance, Navigation, and Control (GNC) algorithms were rigorously validated through Model-In-The-Loop (MIL) simulations. A custom 6DOF trajectory simulator generated artificial sensor data to verify the convergence of the Error-State Kalman Filter (ESKF). The closed-loop guidance and control algorithms demonstrated apogee height errors of less than five metres under varied simulated conditions. Furthermore, the non-volatile storage subsystem was subjected to maximum polling frequency constraints. This validated steady-state write latencies, confirming that the flight computer can sustain maximum operational data throughput without experiencing data corruption, buffer overflows, or real-time operating system (RTOS) task starvation.

\newpage
% Con esto le dices a LaTeX: "Muestra todo lo que haya en el .bib aunque no esté citado"
\nocite{*} 

% Cambiamos el título a "Bibliography" (que se usa cuando son lecturas de fondo, no citas directas)
\printbibliography[heading=bibintoc, title={Bibliography}]


\newpage
% =========================================================================
% APPENDICES (No cuentan para el límite de 25 páginas)
% =========================================================================
\appendix
\section{Risk Matrix}
% Requisito: Matriz de riesgos en formato tabla.

% =========================================================================
% CUSTOM COLOURS FOR RISK LEVEL MATRIX 
% =========================================================================
\definecolor{RiskGreen}{HTML}{19B056}   % Low Risk / Level 1-2
\definecolor{RiskYellow}{HTML}{FFD600}  % Medium Risk / Level 3
\definecolor{RiskOrange}{HTML}{F07D18}  % High Risk / Level 4
\definecolor{RiskRed}{HTML}{C00000}     % Critical Risk / Level 5

The underlying engineering rationale and methodology employed to populate the subsystem risk matrix are detailed below:

% =========================================================================
% TABLE 1: LIKELIHOOD CRITERIA WITH COLOURS (Cabe en una sola página)
% =========================================================================
\begin{table}[H]
    \centering
    \caption{Risk Assessment: Likelihood Criteria and Technical Metrics.}
    \label{tab:likelihood_criteria}
    \vspace{0.2cm}
    \begin{tabularx}{\textwidth}{|>{\centering\arraybackslash}m{1.5cm}|>{\centering\arraybackslash}m{2.5cm}|X|}
        \hline
        \rowcolor{EtseNavy}
        \textbf{\color{white}Level} & \textbf{\color{white}Descriptor} & \textbf{\color{white}Technical Criteria / Estimated Frequency} \\ \hline
        \cellcolor{RiskGreen}1 & \textbf{Remote} & Extremely unlikely to occur. Theoretical probability is negligible; never observed in similar student or experimental rocketry projects. \\ \hline
        \cellcolor{RiskGreen}2 & \textbf{Unlikely} & Rare but possible. Has occurred isolatedly in the association's past launch history or in other university rocketry teams. \\ \hline
        \cellcolor{RiskYellow}3 & \textbf{Possible} & Occasional occurrence expected in unproven SRAD systems. Likely to happen at least once during ground static fire campaigns or flight if environmental variables fluctuate. \\ \hline
        \cellcolor{RiskOrange}4 & \textbf{Likely} & High probability of occurrence. Repeated defects in legacy components or complex algorithms without full qualification. Will occur without robust mitigation. \\ \hline
        \cellcolor{RiskRed}5 & \textbf{Frequent} & Systematic or inherent to the current design architecture. Expected to occur repeatedly during development phases until critical redesigns are implemented. \\ \hline
    \end{tabularx}
\end{table}

% =========================================================================
% TABLE 2: SEVERITY CRITERIA WITH COLOURS (Cabe en una sola página)
% =========================================================================
\begin{table}[H]
    \centering
    \caption{Risk Assessment: Severity Criteria and Mission Impact.}
    \label{tab:severity_criteria}
    \vspace{0.2cm}
    \begin{tabularx}{\textwidth}{|>{\centering\arraybackslash}m{1.5cm}|>{\centering\arraybackslash}m{2.5cm}|X|}
        \hline
        \rowcolor{EtseNavy}
        \textbf{\color{white}Level} & \textbf{\color{white}Descriptor} & \textbf{\color{white}Effect on Mission, Vehicle, and Safety} \\ \hline
        \cellcolor{RiskGreen}1 & \textbf{Negligible} & Zero or minimal degradation of mission performance. Minor data anomalies that do not compromise operations (e.g., noise in a secondary sensor). Mission success remains at 100\%. \\ \hline
        \cellcolor{RiskGreen}2 & \textbf{Minor} & Light performance degradation or partial loss of non-critical subsystems (e.g., temporary telemetry packet drop). The rocket reaches a safe apogee and primary objectives are achieved. \\ \hline
        \cellcolor{RiskYellow}3 & \textbf{Major} & Failure to achieve primary mission objectives (e.g., airbrakes failing to deploy, preventing active apogee control, or payload data corruption). The rocket remains structurally intact and is recovered safely. \\ \hline
        \cellcolor{RiskOrange}4 & \textbf{Critical} & Absolute mission failure resulting in total loss or severe destruction of the NAOS EVO launcher (e.g., ballistic descent due to recovery ejection failure, structural breakup during ascent). No threat to human safety. \\ \hline
        \cellcolor{RiskRed}5 & \textbf{Catastrophic} & Direct risk of severe injuries or fatalities to team members, EuRoC judges, or the general public. Uncontrolled explosion on the launch rail or severe destruction of ground pad infrastructure. \\ \hline
    \end{tabularx}
\end{table}

\subsection{Risk Register and Mitigation Matrix}
The quantified risks across the different vehicle subsystems, ranging from aerostructures to propulsion and software architectures, are detailed in Table~\ref{tab:risk_register}. Appropriate mitigation actions have been actively implemented to reduce all residual risks to acceptable engineering levels.

% =========================================================================
% TABLE 3: RISK REGISTER WITH EXACT CONDITIONAL COLOURS (MULTIPAGE FIX)
% =========================================================================
\small
\begin{xltabular}{\textwidth}{|>{\centering\arraybackslash}p{1.0cm}|>{\raggedright\arraybackslash}p{2.3cm}|>{\centering\arraybackslash}p{0.5cm}|>{\centering\arraybackslash}p{0.5cm}|>{\centering\arraybackslash}p{0.6cm}|X|>{\centering\arraybackslash}p{1.4cm}|}
    
    % Encabezado Primera Página
    \caption{NAOS EVO Comprehensive Subsystem Risk Register.} \label{tab:risk_register} \\ \hline
    \rowcolor{EtseNavy}
    \textbf{\color{white}ID} & \textbf{\color{white}Risk Title} & \textbf{\color{white}L} & \textbf{\color{white}S} & \textbf{\color{white}RI} & \textbf{\color{white}Mitigation Action} & \textbf{\color{white}Residual} \\ \hline
    \endfirsthead

    % Encabezado Páginas Siguientes
    \multicolumn{7}{c}{{\bfseries \tablename\ \thetable{} -- Continued from previous page}} \\ \hline
    \rowcolor{EtseNavy}
    \textbf{\color{white}ID} & \textbf{\color{white}Risk Title} & \textbf{\color{white}L} & \textbf{\color{white}S} & \textbf{\color{white}RI} & \textbf{\color{white}Mitigation Action} & \textbf{\color{white}Residual} \\ \hline
    \endhead

    % Pie de tabla intermedio
    \hline
    \multicolumn{7}{|r|}{{Continued on next page...}} \\ \hline
    \endfoot

    % Pie de tabla final
    \hline
    \endlastfoot
    
    R-001 & Mass growth & 2 & 1 & \cellcolor{RiskGreen}2 & Apply safety margins; run sensitivity simulations to assess mass variation impact on performance; focus on mass budget accuracy. & Low \\ \hline
    
    R-002 & Fins misalignment & 2 & 3 & \cellcolor{RiskYellow}6 & Run sensitivity simulations considering independent fin misalignment angles via Monte Carlo; verify post-manufacturing tolerances. & Medium \\ \hline
    
    R-003 & Fuselage misalignment & 2 & 3 & \cellcolor{RiskYellow}6 & Redundant antennas; use a high-gain ground station setup to compensate for radiation pattern distortion. & Low \\ \hline
    
    R-004 & Telemetry loss & 2 & 1 & \cellcolor{RiskGreen}2 & Implement robust link budgets; execute pre-flight RF path testing. & Low \\ \hline
    
    R-005 & Parachute /drogue breaking & 2 & 4 & \cellcolor{RiskOrange}8 & Structural verification of deployment bags, shock cords, and riser lines; implement high safety factors. & Low \\ \hline
    
    R-006 & No Parachute Deployment & 2 & 5 & \cellcolor{RiskRed}10 & Redundant launcher state estimation algorithms (GNSS, barometer, Kalman filter); extensive OBC signal verification. & Low \\ \hline
    
    R-007 & No Drogue Deployment & 2 & 5 & \cellcolor{RiskRed}10 & Subsystem power supply redundancy; ground validation of ejection forces; redundant pyrotechnic igniters. & Medium \\ \hline
    
    R-008 & Engine ignition failure & 2 & 5 & \cellcolor{RiskRed}10 & Rigorous pre-flight igniter characterisation; extensive electrical continuity testing protocol. & Low \\ \hline
    
    R-009 & GNC instability & 2 & 4 & \cellcolor{RiskOrange}8 & Execute Hardware-in-the-Loop (HIL) testing; conduct comprehensive 6-DOF flight dynamics simulations. & Low \\ \hline
    
    R-010 & Delay in long-lead items & 3 & 3 & \cellcolor{RiskYellow}9 & Early procurement ordering and planning; active identification of alternative backup suppliers. & Medium \\ \hline
    
    R-011 & Nozzle throat blockage & 3 & 3 & \cellcolor{RiskYellow}9 & Increase residence time by managing chamber pressure; reduce Al percentage in the APCP composition; validate via static fire tests. & Medium \\ \hline
    
    R-012 & Burning rate deviation & 2 & 2 & \cellcolor{RiskGreen}4 & Extensive propellant characterisation through open-air and subscale static fire tests to verify $a$ and $n$ coefficients. & Low \\ \hline
    
    R-013 & Flight controller CPU lock & 1 & 4 & \cellcolor{RiskGreen}4 & Implement Direct Memory Access (DMA) for long blocking tasks; Software-in-the-Loop (SIL) and HIL validation. & Low \\ \hline
    
    R-014 & Kalman filter non-convergence & 2 & 3 & \cellcolor{RiskYellow}6 & Fallback to raw barometric data if covariance bounds exceed predetermined thresholds; extensive HIL testing. & Medium \\ \hline
    
    R-015 & Ejection system gas leak & 2 & 5 & \cellcolor{RiskRed}10 & Ground deployment testing of CO$_2$ ejection charges; improve bulkhead sealing interfaces and O-ring tolerances. & Low \\ \hline
    
    R-016 & Parachute lines tangling & 2 & 4 & \cellcolor{RiskOrange}8 & Develop and strictly follow a standardised, documented packing checklist and deployment line sleeves. & Low \\ \hline
    
    R-017 & Premature drogue deployment & 1 & 5 & \cellcolor{RiskGreen}5 & Correctly size and select the mechanical retention spring; complete dynamic load analysis of the retention mechanism. & Low \\ \hline
    
    R-018 & Airbrakes mechanism jam & 2 & 2 & \cellcolor{RiskGreen}4 & Ground load testing of the actuation mechanism; guarantee low-friction linear guide tolerances at the deployment rails. & Low \\ \hline
    
    R-019 & GS Antenna misalignment & 2 & 3 & \cellcolor{RiskYellow}6 & Pre-flight calibration of the automatic antenna tracker; implement manual override controls at the ground station. & Low \\ \hline
    
    R-020 & GS Software (GUI) crash & 2 & 2 & \cellcolor{RiskGreen}4 & Software stress testing using simulated noisy data streams; implement automated quick-restart scripts. & Low \\ \hline
\end{xltabular}

\newpage
\section{Vehicle Development Status}
% Requisito: Fotografías y/o enlaces a vídeos del estado actual.

This section outlines the current development and manufacturing status of NAOS EVO graphical evidence. Please note that complementary video footage is provided as supplementary material in parallel with this progress report.

\begin{figure}[H] % <-- Modificado a [H] para forzar la posición exacta
    \centering
    \includegraphics[width=0.6\textwidth]{Pictures/APCP Grains} 
    \caption{Solid rocket motor APCP grain production and quality inspection status.}
    \label{fig:apcp_grains}
\end{figure}

\begin{figure}[H] % <-- Modificado a [H] para forzar la posición exacta
    \centering
    \includegraphics[width=0.6\textwidth]{Pictures/Ground Station UI} 
    \caption{Ground Station GUI}
    \label{fig:recovery_test}
\end{figure}

% =========================================================================
% FIGURE A.1: CÁMARA A BORDO 4
% =========================================================================
\begin{figure}[H]
    \centering
    \includegraphics[width=0.7\textwidth]{Pictures/Pictures_APENDIX/Camara_a_bordo_4.jpg}
    \caption{Onboard flight camera frame.}
    \label{fig:onboard_cam_4}
\end{figure}

% =========================================================================
% FIGURE A.2: COHETE NAOS DESPEGUE
% =========================================================================
\begin{figure}[H]
    \centering
    \includegraphics[width=0.7\textwidth]{Pictures/Pictures_APENDIX/Cohete_NAOS_Despegue.jpeg}
    \caption{NAOS EVO launcher liftoff.}
    \label{fig:naos_liftoff}
\end{figure}

% =========================================================================
% FIGURE A.3: COHETE PROMOCIONAL 2 RED
% =========================================================================
\begin{figure}[H]
    \centering
    \includegraphics[width=0.7\textwidth]{Pictures/Pictures_APENDIX/Cohete_Promocional_2_reducido.jpg}
    \caption{NAOS EVO Promotional Picture.}
    \label{fig:naos_promo_2}
\end{figure}

% =========================================================================
% FIGURE A.4: PARACAIDAS LANZAMIENTO 2
% =========================================================================
\begin{figure}[H]
    \centering
    \includegraphics[width=0.7\textwidth]{Pictures/Pictures_APENDIX/Paracaidas_lanzamiento_2.jpg}
    \caption{Old version Main Parachute inflation canopy.}
    \label{fig:parachute_deployment_2}
\end{figure}

% =========================================================================
% FIGURE A.5: VUELO 3
% =========================================================================
\begin{figure}[H]
    \centering
    \includegraphics[width=0.7\textwidth]{Pictures/Pictures_APENDIX/Vuelo_3.jpg}
    \caption{NAOS old version powered launchpad clearance.}
    \label{fig:flight_tracking_3}
\end{figure}


\section{Operational Checklists}
% Requisito: Procedimientos paso a paso.

% =========================================================================
% LAUNCH PAD OPERATIONS CHECKLIST SHEET
% =========================================================================
\subsection{Launch Pad and Countdown Operations Protocol}

This protocol defines the sequential operations executed at the launch pad and the Launch Control Centre (LCC) during the definitive countdown window. The estimated timeline for completion is 1 hour and 10 minutes.

\begin{table}[H]
    \centering
    \small
    \caption{Launch Pad Operational Metadata and Logistics.}
    \label{tab:pad_metadata}
    \vspace{0.2cm}
    \begin{tabularx}{\textwidth}{|>{\cellcolor{LightGrayTable}\bfseries\color{EtseNavy}}l|X|>{\cellcolor{LightGrayTable}\bfseries\color{EtseNavy}}l|X|}
        \hline
        Estimated Time & 1 h 10 min & Target Vehicle & NAOS EVO \\ \hline
        Subsystem Link & Propulsion / Avionics Integration & Control Authority & Launch Control Centre (LCC) \\ \hline
    \end{tabularx}
\end{table}

% Definición del estilo de columnas para xltabular
\footnotesize
\begin{xltabular}{\textwidth}{|>{\centering\arraybackslash}p{0.6cm}|X|X|>{\centering\arraybackslash}p{2.8cm}|>{\centering\arraybackslash}p{1.0cm}|}
    
    % Encabezado de la primera página
    \caption{Sequential Launch Pad and Countdown Verification Matrix.} \label{tab:launch_pad_checklist} \\ \hline
    \rowcolor{EtseNavy}
    \textbf{\color{white}Step} & \textbf{\color{white}Task Description} & \textbf{\color{white}Technical / Operational Notes} & \textbf{\color{white}Location} & \textbf{\color{white}Status} \\ \hline
    \endfirsthead

    % Encabezado que se repetirá en las páginas siguientes
    \multicolumn{5}{c}{{\bfseries \tablename\ \thetable{} -- Continued from previous page}} \\ \hline
    \rowcolor{EtseNavy}
    \textbf{\color{white}Step} & \textbf{\color{white}Task Description} & \textbf{\color{white}Technical / Operational Notes} & \textbf{\color{white}Location} & \textbf{\color{white}Status} \\ \hline
    \endhead

    % Pie de tabla en páginas intermedias
    \hline
    \multicolumn{5}{|r|}{{Continued on next page...}} \\ \hline
    \endfoot

    % Pie de tabla definitivo al final del todo
    \hline
    \endlastfoot
    
    % --- AUDIOVISUAL REPORT ---
    \multicolumn{5}{|l|}{\cellcolor{LightGrayTable}\bfseries\color{EtseNavy}1. Audiovisual Documentation} \\ \hline
    1 & Initiate operational media recording campaign. & Coordinate camera positions across public, pad, and control segments (if possible). & Pad / LCC / Public & \\ \hline
    
    % --- ROCKET OFF THE RAIL ---
    \multicolumn{5}{|l|}{\cellcolor{LightGrayTable}\bfseries\color{EtseNavy}2. Pre-Integration Verification (Vehicle off the Launch Rail)} \\ \hline
    2 & Verify motor ignition safety interlocks are open. & Ensure open-circuit state prior to structural integration. & Launch Pad & \\ \hline
    3 & Verify drogue recovery ejection safety interlocks are open. & Ensure open-circuit state across primary and secondary lines. & Launch Pad & \\ \hline
    4 & Verify tender-descender ejection interlocks are open. & Ensure open-circuit state across the secondary recovery bulkhead. & Launch Pad & \\ \hline
    5 & Confirm "Remove Before Flight" (RBF) tag is engaged on the motor igniter electronics. & Mechanical isolation loop must remain actively installed. & Launch Pad & \\ \hline
    6 & Confirm RBF tag is engaged on the drogue ejection controller. & Mechanical isolation loop must remain actively installed. & Launch Pad & \\ \hline
    7 & Confirm RBF tag is engaged on the tender-descender parachute electronics. & Mechanical isolation loop must remain actively installed. & Launch Pad & \\ \hline
    8 & Switch the primary avionics battery power bus to ON. & Align with Avionics Checklist Step 2 (system boot). & Launch Pad & \\ \hline
    9 &Verify live RF telemetry link status continuity. & Confirm data parsing frequency at the ground station terminal. & Launch Pad / LCC & \\ \hline 
    10 & Execute dedicated avionics stack diagnostic checklist. & Reference internal PCB diagnostic matrices. & Launch Pad & \\ \hline
    11 &Validate GNSS telemetry string reception at ground tracking stations. & Align with Avionics Checklist Step 15 ($+4$ satellites). & Launch Pad / LCC & \\ \hline
    
    % --- VIDEO SYSTEMS ---
    \multicolumn{5}{|l|}{\cellcolor{LightGrayTable}\bfseries\color{EtseNavy}3. Video and Recording Systems Setup} \\ \hline
    12 & Position and align internal/external onboard Video Recording units. & Verify field of view parameters. & Pad / LCC / Public & \\ \hline
    13 & Mount and lock auxiliary ground recording cameras. & Position for structural and trajectory tracking (if possible). & Pad / LCC / Public & \\ \hline
    
    % --- IGNITION SYSTEM ---
    \multicolumn{5}{|l|}{\cellcolor{LightGrayTable}\bfseries\color{EtseNavy}4. Ignition Subsystem Architecture} \\ \hline
    14 & Launch Pad crew issues operational status update to LCC. & Notification of ignition system preparation commencement. & Launch Pad & \\ \hline
    15 & [TBC] Connect ignition umbilical cable to the physical ignition control box. & Ensure secure terminal block seating. & LCC & \\ \hline
    \rowcolor{RiskYellow}
    16 & [TBC] Extract and safely store the physical key from the ignition control box. & Critical safety lockout protocol. & LCC & \\ \hline
    \rowcolor{RiskYellow}
    17 & Confirm the ignition control box arming switch is physically set to OFF. & Visual verification required. & LCC & \\ \hline
    18 & Launch Pad crew transmits an operational update to LCC. & Estimated $T-30\,\text{minutes}$ until launch window opens. & Launch Pad & \\ \hline
    \rowcolor{RiskOrange}
    19 & Issue official public and range warning notice. & Formal $T-30\,\text{minutes}$ audible and radio broadcast. & LCC & \\ \hline
    
    % --- LAUNCH RAIL OPERATIONS ---
    \multicolumn{5}{|l|}{\cellcolor{LightGrayTable}\bfseries\color{EtseNavy}5. Launch Rail Integration (LAUNCHPAD)} \\ \hline
    20 & Integrate the NAOS EVO vehicle into the launch rail guides. & Smooth transition check across structural rail buttons. & Launch Pad & \\ \hline
    21 & Adjust the launch rail elevation and azimuth angles. & Match nominal trajectory simulation coordinates.  Consider Wind Conditions& Launch Pad & \\ \hline
    22 & [TBC] Protect composite aerodynamic fin units with sun thermal protector. & Deploy if ambient solar radiation limits are exceeded. & Launch Pad & \\ \hline
    23 & Take picture as-built vehicle integration state on the rail. & High-resolution photographic reference capture (if possible). & Launch Pad & \\ \hline
    24 & Insert the pyrotechnic igniter into the motor throat. & Verify depth alignment matches the nozzle throat gauge marking. & Launch Pad & \\ \hline
    25 & Apply thermal-resistant adhesive tape to the nozzle exit plane. & Temporary retention matrix for structural cable support. & Launch Pad & \\ \hline
    26 & [TBC] Connect umbilical leads from the pad ignition box to the igniter. & Verify electrical terminal connection integrity. & Launch Pad & \\ \hline
    27 & Initiate launchpad active recording sequence. & Verification via internal wireless link or physical prompt (if possible). & Launch Pad & \\ \hline
    28 & Initiate active recording on the fixed public-facing camera tracking. & Confirm operational status (if possible). & Launch Pad & \\ \hline
    29 & Initiate active recording on the fixed launch-button camera tracking. & Continuous control panel monitoring (if possible). & LCC & \\ \hline
    30 & Initiate active recording on Public Segment. & Multi-angle videorecording public. & Public Segment & \\ \hline
    \rowcolor{RiskYellow}
    31 & Activate internal onboard high-definition camera. & Target resolution: $1080\,\text{p}$ at $60\,\text{fps}$. & Launch Pad & \\ \hline
    \rowcolor{RiskYellow}
    32 & Confirm all operational personnel have evacuated to safe-zones. & Absolute clear perimeter validation. & Launch Pad & \\ \hline
    \rowcolor{RiskYellow}
    33 & Verify the designated Ignition Officer holds the physical arming key. & Verbal and visual validation required. & LCC & \\ \hline
    34 & Transmit Launch Pad clearance confirmation to the LCC. & Ready for system energisation phase. & Launch Pad & \\ \hline
    \rowcolor{RiskGreen}
    35 & Obtain formal launch authorisation from the Ignition Officer at Launchpad. & - & Launch Pad & \\ \hline
    \rowcolor{RiskYellow}
    36 & Broadcast imminent coupling warning of the ignition system lines. & Radio notification issued by the Ignition Officer. & Launch Pad & \\ \hline
    \rowcolor{RiskRed}
    37 & Connect the pyrotechnic igniter leads to the main firing line. & Final hazardous electrical connection step. & Launch Pad & \\ \hline
    \rowcolor{RiskGreen}
    38 & Obtain formal clearance confirmation from the Media Director. & Ground and tracking systems fully active (if possible). & Pad / LCC / Public & \\ \hline
    \rowcolor{RiskGreen}
    39 & Obtain formal telemetry clearance from GNSS ground stations. & Signal lock and data logging integrity verified. & LCC & \\ \hline
    \rowcolor{RiskGreen}
    40 & Obtain formal structural clearance from Launch Pad Director. & Rail and structural vectors verified. & Launch Pad & \\ \hline
    \rowcolor{RiskGreen}
    41 & Execute airspace safety sweep for unauthorized aerial vehicles. & Confirm zero traffic within the designated launch corridor. & LCC & \\ \hline
    \rowcolor{RiskGreen}
    42 & Initiate definitive countdown sequence. & Final terminal count sequence engagement. & LCC & \\ \hline
    
    % --- POST-RECOVERY ---
    \multicolumn{5}{|l|}{\cellcolor{LightGrayTable}\bfseries\color{EtseNavy}6. Post-Recovery Procedures} \\ \hline
    43 & Safe the vehicle and isolate the internal Li-Ion battery packs. & Extract packs to prevent residual discharge or hazard. & Recovery Site & \\ \hline
    44 & Terminate onboard camera recording arrays and extract data. & Secure video logs for post-flight telemetry analysis. & Recovery Site & \\
\end{xltabular}

\subsection{Avionic Bay/Ground Station Systems Checklist}

This section details the pre-flight verification protocol designed to ensure the operational readiness of the NAOS EVO avionics bay and corresponding ground station infrastructure.

% =========================================================================
% TABLE 1: METADATA & HEADER BLOCK
% =========================================================================
\begin{table}[H]
    \centering
    \small
    \caption{Avionics Bay / Ground Station Systems Checklist Metadata.}
    \label{tab:checklist_metadata}
    \vspace{0.2cm}
    \begin{tabularx}{\textwidth}{|>{\cellcolor{LightGrayTable}\bfseries\color{EtseNavy}}l|X|>{\cellcolor{LightGrayTable}\bfseries\color{EtseNavy}}l|X|}
        \hline
        Date & 06/07/2026 & Time & \\ \hline
        Member(s) & \multicolumn{3}{l|}{José Reyes Espina García, Alejandro Ariza} \\ \hline
        Approved by & \multicolumn{3}{l|}{} \\ \hline
        \hline
        \rowcolor{LightGrayTable}
        \multicolumn{4}{|p{\dimexpr\textwidth-2\tabcolsep-2\arrayrulewidth}|}{\textbf{\color{EtseNavy}Brief Description:}} \\ \hline
        \multicolumn{4}{|X|}{Pre-flight checklist protocol to verify structural and electrical integration constraints of the NAOS EVO rocket avionics bay and the associated ground station.} \\ \hline
    \end{tabularx}
\end{table}

% =========================================================================
% TABLE 2: LOGISTICS, TOOLS AND EQUIPMENT (CORREGIDO)
% =========================================================================
\begin{table}[H]
    \centering
    \small
    \caption{Required Field Tools, Protective Equipment, and Subsystem Components.}
    \label{tab:checklist_tools}
    \vspace{0.2cm}
    \begin{tabularx}{\textwidth}{|X|X|X|}
        \hline
        \rowcolor{EtseNavy}
        \textbf{\color{white}Tools} & \textbf{\color{white}Protective Equipment} & \textbf{\color{white}Vehicle Parts} \\ \hline
        \begin{itemize}
            \item Allen key kit
            \item Screwdrivers
            \item M3 screws
            \item Pliers
            \item Digital multimeter
        \end{itemize} & 
        \begin{itemize}
            \item Safety gloves
            \item Protective glasses
        \end{itemize} & 
        \begin{itemize}
            \item Custom PCBs
            \item Ground Station unit
            \item GS Tracking Antennas
            \item Avionics Bay structural frame
            \item Lithium Ions battery pack
        \end{itemize} \\ \hline
    \end{tabularx}
\end{table}

% =========================================================================
% TABLE 3: CORE OPERATIONAL CHECKLIST STEPS (UNIFIED)
% =========================================================================
\begin{table}[H]
    \centering
    \footnotesize
    \caption{Sequential Pre-Flight Integration and Verification Protocol.}
    \label{tab:core_checklist_steps}
    \vspace{0.2cm}
    \begin{tabularx}{\textwidth}{|>{\centering\arraybackslash}p{0.6cm}|X|X|>{\centering\arraybackslash}p{0.9cm}|>{\centering\arraybackslash}p{1.2cm}|>{\centering\arraybackslash}p{0.9cm}|}
        \hline
        \rowcolor{EtseNavy}
        \textbf{\color{white}Step} & \textbf{\color{white}Task Description} & \textbf{\color{white}Additional Comments} & \textbf{\color{white}Status} & \textbf{\color{white}Approved} & \textbf{\color{white}Time} \\ \hline
        
        \rowcolor{LightGrayTable}
        1 & Test all internal batteries are charged up to 100\% ($4.2\,\text{V}$). & Execute verification prior to structural mounting within the upper avionics bay rack. & & & 2 min \\ \hline
        
        2 & Verify arming switches operate nominally to wake up onboard electronic systems. & ONLY remove Baseboard on/off, Cats Vega on/off, and Active Control on/off interfaces sequentially. & & & 2 min \\ \hline
        
        \rowcolor{LightGrayTable}
        3 & Verify Powerboard $5\,\text{V}$ and $3.3\,\text{V}$ electrical bus outputs. & Measure lines once all battery packs are fully assembled and verified. & & & 2 min \\ \hline
        
        4 & Characterise Cats Vega nominal operational status. & Confirm automated audible buzzing/beeping response sequence. & & & 2 min \\ \hline
        
        \rowcolor{LightGrayTable}
        5 & Verify Powerboard PWM output functionality during active control system deployment. & Conduct a sweep tracking from 0\% to 100\% mechanical petal deflection. & & & 5 min \\ \hline
        
        6 & Confirm Powerboard safety relays function nominally. & Audio verification of mechanical relay actuation click. & & & 2 min \\ \hline
        
        \rowcolor{LightGrayTable}
        7 & Monitor Baseboard finite state machine status transitions (IDLE). & System must generate a single audible beep tone. & & & 2 min \\ \hline
        
        8 & Validate telemetry link continuity in IDLE state. & Requires explicit confirmation loop from the ground station terminal. & & & 2 min \\ \hline
        
        \rowcolor{LightGrayTable}
        9 & Initiate state machine event change sequence to CALIBRATION. & Requires telemetry confirmation from the ground station operator. & & & 1 min \\ \hline
        
        10 & Verify Baseboard sensor clusters read within expected analytical bounds. & Verified if the baseboard beeps, telemetry reports data, and the SD logging card is correctly locked. & & & 2 min \\ \hline
        
        \rowcolor{LightGrayTable}
        11 & Confirm onboard Baseboard data logging pipelines function nominally. & Verify data writing loop by extracting and parsing sample files from the SD micro-bus. & & & 5 min \\ \hline
        
        12 & Extract physical mechanical arming pins from the pyrotechnic switch matrix. & Verify a count of exactly 4 arming pins are cleared (On/off-PYRO). & & & 2 min \\ \hline
        
        \rowcolor{LightGrayTable}
        13 & Re-verify comprehensive data telemetry link margins. & Explicit link budget confirmation from the Ground Station hub. & & & 4 min \\ \hline
        
        14 & Confirm system has achieved absolute calibration state status. & Visual assessment of telemetry data stream. & & & 1 min \\ \hline
        
        \rowcolor{LightGrayTable}
        15 & Validate operational GNSS satellite communication tracking. & A minimum threshold of $+4$ locked satellites is strictly required. & & & 2 min \\ \hline
        
        16 & Perform final telemetry sanity checks on sensor cluster readings. & Final verbal confirmation of data integrity to be granted by ground station director. & & & 2 min \\ \hline
    \end{tabularx}
\end{table}

% =========================================================================
% REMARKS & PROTOCOL NOTES
% =========================================================================
\noindent
\begin{minipage}{\textwidth}
\small
\textbf{\color{EtseNavy}Operational Implementation Remarks:}
\begin{itemize}
    \item \textbf{Steps 1--11:} Must be completed sequentially during vehicle integration operations inside the preparation tent.
    \item \textbf{Steps 12--13:} Executed exclusively once the launch vehicle is safely integrated on the launch rail.
    \item \textbf{Steps 14--16:} To be performed only upon receiving a formal launch countdown window confirmation from the EuRoC Launch Director.
\end{itemize}
\end{minipage}

\vspace{0.4cm}

% =========================================================================
% TABLE 4: SIGN-OFF & FIELD VALIDATION SIGNATURES
% =========================================================================
\begin{table}[H]
    \centering
    \small
    \caption{Field Sign-Off and Formal Authorization Blocks.}
    \label{tab:checklist_signatures}
    \vspace{0.2cm}
    \begin{tabularx}{\textwidth}{|X|X|X|}
        \hline
        \multicolumn{3}{|p{\dimexpr\textwidth-2\tabcolsep-2\arrayrulewidth}|}{We hereby confirm that all tasks of this checklist were completed as guided in the Assembly Instruction Sheet Ref. \hrulefill} \\ \hline
        \rule{0pt}{1.5cm} Signature of Member 1: & Signature of Member 2: & Signature of Member 3: \\ 
        \hrulefill & \hrulefill & \hrulefill \\ \hline
        \rule{0pt}{1.5cm} Responsible Signature: & Responsible Signature: & Responsible Signature: \\ 
        \hrulefill & \hrulefill & \hrulefill \\ \hline
    \end{tabularx}
\end{table}

\end{document}